\documentclass[11pt,a4paper]{article}

\usepackage[utf8]{inputenc}
\usepackage[T1]{fontenc}
\usepackage{amsmath}
\usepackage{geometry}
\usepackage{hyperref}
\usepackage{verbatim}

\geometry{margin=2.5cm}

\newenvironment{codeblock}{\verbatim}{\endverbatim}

\title{Guía de uso del evaluador de expresiones}
\author{Luciano David Duarte}
\date{ALP 2025 - Licenciatura en Ciencias de la Computación}

\begin{document}

\maketitle

\section{Objetivo de la guía}

Esta guía explica cómo compilar, ejecutar y probar el evaluador de expresiones
matemáticas desarrollado para el trabajo final. El programa permite escribir
expresiones en una variable principal \texttt{x}, evaluarlas en un punto y
calcular su derivada usando números duales.

El proyecto también incluye procesamiento de archivos, manejo explícito de
errores, optimizaciones conservadoras y una API interna que permite evaluar con
entornos de variables.

\section{Requisitos}

Para usar el proyecto se necesita tener instalado:

\begin{itemize}
  \item GHC compatible con el proyecto;
  \item Cabal;
  \item una terminal de Windows, PowerShell o WSL.
\end{itemize}

Los comandos principales de esta guía están escritos para PowerShell desde la
raíz del repositorio:

\begin{codeblock}
cd C:\Users\lucho\ALP-2025-LCC
\end{codeblock}

\section{Compilación}

Para compilar biblioteca, ejecutable y tests:

\begin{codeblock}
cabal build all
\end{codeblock}

Si Cabal responde \texttt{Up to date}, significa que el proyecto ya está
compilado con los cambios actuales.

\section{Modo interactivo}

El modo interactivo permite ingresar una expresión, pedir un valor para
\texttt{x} y obtener el valor de la función junto con su derivada.

\begin{codeblock}
cabal run ALP2025-LCC
\end{codeblock}

Ejemplo de uso:

\begin{codeblock}
>>> sin(x) + x^2
>> Ingrese el valor de x: 1
f(1) = 1.8414709848078965
f'(1) = 2.5403023058681398
\end{codeblock}

Para salir del modo interactivo se puede escribir:

\begin{codeblock}
salir
quit
exit
\end{codeblock}

\section{Procesamiento de archivos}

El programa puede leer varias expresiones desde un archivo. Cada línea útil debe
tener este formato:

\begin{codeblock}
expresion @ valor_de_x
\end{codeblock}

Ejemplo:

\begin{codeblock}
x^2 + 2*x + 1 @ 3
sin(x) @ pi/2
log(x) + pi @ e
\end{codeblock}

Para procesar un archivo:

\begin{codeblock}
cabal run ALP2025-LCC -- -f examples\basico.txt
\end{codeblock}

También se puede usar la opción larga:

\begin{codeblock}
cabal run ALP2025-LCC -- --file examples\basico.txt
\end{codeblock}

En algunas instalaciones de Windows puede bloquearse el ejecutable generado por
Cabal. En ese caso se puede correr el mismo flujo con \texttt{runghc}:

\begin{codeblock}
cabal exec -- runghc -isrc -iapp app\Main.hs -f examples\defensa.txt
\end{codeblock}

\section{Comentarios en archivos}

Los archivos aceptan comentarios de línea:

\begin{codeblock}
-- Esto es un comentario
x^2 @ 3 -- comentario al final de la línea
\end{codeblock}

También aceptan comentarios de bloque:

\begin{codeblock}
{- Comentario
   de varias líneas -}

(x + {- comentario interno -} 1) * 2 @ 3
\end{codeblock}

\section{Sintaxis soportada}

\subsection{Operadores}

\begin{center}
\begin{tabular}{ll}
\textbf{Operador} & \textbf{Significado} \\
\hline
\texttt{+} & suma \\
\texttt{-} & resta y menos unario \\
\texttt{*} & multiplicación \\
\texttt{/} & división \\
\texttt{\^{}} & potencia \\
\end{tabular}
\end{center}

La precedencia usada por el parser es:

\begin{codeblock}
potencia > menos unario > multiplicación/división > suma/resta
\end{codeblock}

Por ejemplo, \texttt{-x\^{}2} se interpreta como \texttt{-(x\^{}2)}.

\subsection{Funciones}

\begin{center}
\begin{tabular}{ll}
\textbf{Función} & \textbf{Dominio principal} \\
\hline
\texttt{sin(x)}, \texttt{cos(x)} & todo número real \\
\texttt{tan(x)} & donde \texttt{cos(x)} no sea cero \\
\texttt{sinh(x)}, \texttt{cosh(x)}, \texttt{tanh(x)} & todo número real \\
\texttt{arsinh(x)} & todo número real \\
\texttt{arcosh(x)} & \texttt{x >= 1} \\
\texttt{artanh(x)} & \texttt{|x| < 1} \\
\texttt{sqrt(x)} & \texttt{x >= 0} \\
\texttt{log(x)} & \texttt{x > 0} \\
\texttt{exp(x)} & todo número real \\
\end{tabular}
\end{center}

\subsection{Constantes}

Se soportan las constantes:

\begin{codeblock}
pi
e
\end{codeblock}

El valor ubicado a la derecha de \texttt{@} puede ser un número o una expresión
constante, por ejemplo:

\begin{codeblock}
sin(x) @ pi/2
log(x) @ e
\end{codeblock}

\section{Optimización}

Antes de evaluar, el programa aplica optimizaciones simples. Algunos ejemplos:

\begin{codeblock}
x + 0  -> x
x * 1  -> x
x / 1  -> x
x^1    -> x
2 + 3  -> 5
\end{codeblock}

Las optimizaciones son conservadoras. Esto significa que no deben ocultar errores
que el evaluador debería detectar. Por eso no se simplifican automáticamente casos
como:

\begin{codeblock}
x / x  -> 1
x * 0  -> 0
x^0    -> 1
1^x    -> 1
\end{codeblock}

Por ejemplo, si \texttt{x = 0}, entonces \texttt{x/x} debe seguir dando división
por cero y \texttt{x\^{}0} debe respetar el caso especial \texttt{0\^{}0}.

\section{Manejo de errores}

Los errores principales se representan de forma tipada y luego se muestran al
usuario con mensajes legibles.

\subsection{División por cero}

\begin{codeblock}
1/0 @ 0
\end{codeblock}

Salida esperada:

\begin{codeblock}
[X] Error al evaluar: Division por cero
\end{codeblock}

\subsection{Error de dominio}

\begin{codeblock}
sqrt(-1) @ 0
\end{codeblock}

Salida esperada:

\begin{codeblock}
[X] Error al evaluar: Error de dominio: sqrt requiere argumento no negativo
\end{codeblock}

\subsection{Variable no definida}

El modo de usuario trabaja con la variable \texttt{x}. Si se usa una variable sin
valor en el entorno, se reporta como error:

\begin{codeblock}
y + 2 @ 1
\end{codeblock}

Salida esperada:

\begin{codeblock}
[X] Error al evaluar: Variable no definida: y
\end{codeblock}

\section{Archivos de ejemplo}

La carpeta \texttt{examples/} incluye casos para probar distintas partes del
lenguaje:

\begin{center}
\begin{tabular}{ll}
\textbf{Archivo} & \textbf{Contenido principal} \\
\hline
\texttt{basico.txt} & expresiones iniciales \\
\texttt{compuestas.txt} & composición de funciones \\
\texttt{constantes.txt} & uso de \texttt{pi} y \texttt{e} \\
\texttt{constantes\_extendido.txt} & constantes en casos más variados \\
\texttt{trigonometricas.txt} & funciones trigonométricas \\
\texttt{hiperbolicas.txt} & funciones hiperbólicas \\
\texttt{todas\_funciones.txt} & combinación de funciones soportadas \\
\texttt{numeros\_negativos.txt} & menos unario y negativos \\
\texttt{precedencia\_operadores.txt} & precedencia y asociatividad \\
\texttt{optimizaciones.txt} & reglas de optimización \\
\texttt{casos\_error.txt} & errores controlados \\
\texttt{test\_comentarios.txt} & comentarios de línea y de bloque \\
\texttt{valores\_especiales.txt} & valores borde simples \\
\texttt{anidamiento\_profundo.txt} & expresiones anidadas \\
\texttt{identidades\_matematicas.txt} & identidades conocidas \\
\texttt{nuevas\_caracteristicas.txt} & funcionalidades agregadas \\
\texttt{defensa.txt} & casos pensados para explicar el trabajo \\
\end{tabular}
\end{center}

\section{Pruebas}

Para correr la suite completa:

\begin{codeblock}
cabal test
\end{codeblock}

La suite actual contiene 125 casos. Incluye una verificación llamada:

\begin{codeblock}
Verificacion de Casos de ejemplo
\end{codeblock}

Esa prueba lee todos los archivos \texttt{.txt} de \texttt{examples/}, verifica
que parseen correctamente y comprueba que la evaluación termine de forma
controlada, ya sea con resultado correcto o con un error esperado.

También se puede correr la suite directamente:

\begin{codeblock}
cabal exec -- runghc -isrc test\TestSuite.hs
\end{codeblock}

\section{Ejemplo recomendado para defensa}

Para mostrar varias decisiones del proyecto en una sola ejecución:

\begin{codeblock}
cabal exec -- runghc -isrc -iapp app\Main.hs -f examples\defensa.txt
\end{codeblock}

Ese archivo muestra:

\begin{itemize}
  \item parsing y precedencia;
  \item derivación automática con números duales;
  \item optimización conservadora;
  \item errores de dominio controlados;
  \item constantes como \texttt{pi} y \texttt{e};
  \item comentarios de bloque dentro de una expresión.
\end{itemize}

\section{Notas finales}

El programa no busca ser un sistema de álgebra computacional completo. Su foco es
representar un lenguaje pequeño de expresiones, interpretarlo de forma segura y
usar números duales para obtener derivadas automáticamente.

\end{document}
